% !TeX root = ../本科微分几何笔记.tex
% 由主文档通过 \input 载入；请编译主文档。
\setcounter{chapter}{9} % 保留原书章号。
\chapter{Gauss--Bonnet 定理与整体曲率}
\label{chap:dg-gb}

\section{边界方向、外转角与局部公式}
给定曲面定向 $\nu$，区域边界的正向规定为：沿边界的单位切向量 $t$ 行进时，$q=\nu\times t$ 指向区域内部。角点的有符号外转角 $\alpha$ 是从入射切向量到出射切向量在定向切平面内的转角；对通常的非退化角点，它等于 $\pi-\beta$，其中 $\beta\in(0,2\pi)$ 为内角。凹角的外转角为负。

\begin{theorem}[坐标圆盘上的 Gauss--Bonnet 公式]
设 $D$ 为包含在一个定向坐标片内、同胚于闭圆盘的紧致曲面区域，边界分段 $C^2$、简单闭合、角点非退化。则
\begin{equation}
 \int_D K\dd A+\int_{\partial D}k_g\dd s+\sum_j\alpha_j=2\pi.
 \label{eq:dg-local-gb}
\end{equation}
曲率积分取正面积，边界按上述正向，$k_g$ 按式~\eqref{eq:dg-darboux-first} 的符号约定。
\label{thm:dg-local-gb}
\end{theorem}
\begin{proof}
在坐标片中对坐标基作 Gram--Schmidt 正交化，得到正向标准正交标架 $e_1,e_2$。边界切向量相对该标架的角记为 $\theta$；由式~\eqref{eq:dg-kg-angle}，各光滑边上有 $k_g\dd s=\dd\theta+\omega$。切向量相对这一标架的总转角加角点跳跃为 $2\pi$：在参数平面上，这是定理~\ref{thm:dg-turning}；从坐标基连续正交化到正交标架的变换在整个圆盘上定义，边界上不增加绕转次数。因此
\[
 \int_{\partial D}k_g\dd s+\sum_j\alpha_j
 =2\pi+\int_{\partial D}\omega.
\]
对一形式 $\omega$ 使用 Green--Stokes 公式，并代入式~\eqref{eq:dg-connection-curvature}，得 $\int_{\partial D}\omega=\int_D\dd\omega=-\int_DK\dd A$，即得结论。此证明依赖前文明确引用的平面转角定理。
\end{proof}
公式把三种转弯合并：区域内部的曲率、边界平滑处的测地转弯以及角点的离散转角。在平面 $K=0$ 时，它退化为转角定理；在测地多边形上，边界测地曲率项为零。

\section{曲面三角形与球面面积}
\begin{corollary}[测地三角形的角盈]
对满足定理~\ref{thm:dg-local-gb} 条件的测地三角形，内角为 $A,B,C$，有
\begin{equation}
 A+B+C-\pi=\int_DK\dd A.
 \label{eq:dg-angle-excess}
\end{equation}
\end{corollary}
\begin{proof}
三条边的 $k_g=0$，三个外角的和为 $3\pi-(A+B+C)$，代入即可。
\end{proof}
对半径 $R$ 球面上的简单测地三角形，只要选择的区域是圆盘且各边为给定的大圆弧，局部公式经细分仍适用，面积为 $R^2(A+B+C-\pi)$。通常采用短大圆弧围成的小三角形以消除区域选择的歧义。单位球面的一个卦限由三条四分之一大圆围成，三个角均为 $\pi/2$，面积为 $\pi/2$。

\begin{example}[球冠的边界曲率核对]
半径 $R$ 球面北部球冠 $0\le\theta\le\theta_0<\pi$ 的面积为 $2\pi R^2(1-\cos\theta_0)$。边界正向是经度增加，测地曲率为 $\cot\theta_0/R$，边长为 $2\pi R\sin\theta_0$，所以
\[
 \int_DK\dd A=2\pi(1-\cos\theta_0),\qquad
 \int_{\partial D}k_g\dd s=2\pi\cos\theta_0.
\]
两项之和为 $2\pi$。经纬度在北极失效不影响曲面公式，可换坐标或将区域细分处理。
\end{example}

\section{欧拉示性数与整体公式}
有限三角剖分将紧致曲面分成有限个三角形，任意两个三角形的交集为空、公共顶点或公共边。设顶点、边、面数分别为 $V,E,F$，定义\emph{欧拉示性数（Euler characteristic）} $\chi=V-E+F$。紧致光滑曲面存在适当三角剖分，且 $\chi$ 与所选剖分无关；这是此处引用的拓扑学前置结果。可细分剖分，使每个小三角形都包含在一个坐标片中，边为分段光滑曲线。

\begin{theorem}[整体 Gauss--Bonnet 定理]
若 $\Sigma$ 是紧致、无边界、可定向光滑曲面，则
\begin{equation}
 \int_\Sigma K\dd A=2\pi\chi(\Sigma).
 \label{eq:dg-global-gb}
\end{equation}
更一般地，对紧致定向曲面区域 $D$，若边界分段光滑且具有非退化角点，则
\begin{equation}
 \int_DK\dd A+\int_{\partial D}k_g\dd s+\sum_j\alpha_j
 =2\pi\chi(D).
 \label{eq:dg-boundary-gb}
\end{equation}
\label{thm:dg-global-gb}
\end{theorem}
\begin{proof}
先考虑无边界情形。对上述细分中的每个三角形应用局部公式，写成 $\int K+\int k_g=\sum\beta-\pi$，其中 $\beta$ 是三个内角。所有公共边出现两次且方向相反，其测地曲率积分抵消；每个顶点周围的内角和为 $2\pi$。于是 $\int_\Sigma K\dd A=2\pi V-\pi F$。闭曲面中每条边属于两个三角形，故 $3F=2E$，从而右端等于 $2\pi(V-E+F)$。

有边界时，记内部和边界顶点数为 $V_i,V_b$，边数为 $E_i,E_b$。内部顶点总角为 $2\pi$，边界顶点总内角为 $\beta_j$，只有内部边积分抵消。因此加上边界外角总和 $\sum_j(\pi-\beta_j)$ 后，左端为 $2\pi V_i+\pi V_b-\pi F$。边界上人为增加的光滑顶点外角为零。利用 $3F=2E_i+E_b$ 及闭合边界多边形的 $E_b=V_b$，所得数恰为 $2\pi(V-E+F)$。证明使用三角剖分存在性，不要求这些分割边本身都是测地线。
\end{proof}

对连通、闭合、可定向曲面，拓扑分类定理给出 $\chi=2-2g$，其中 $g\in\{0,1,2,\ldots\}$ 称为\emph{亏格（genus）}，可直观理解为把手数。该分类定理也作为拓扑学结论引用。于是球面的曲率积分为 $4\pi$，环面为 $0$；对 $g\ge2$ 的曲面，总曲率为负。

\section{应用范围与常见误推}
若紧致无边界连通可定向曲面处处 $K>0$，则 Gauss--Bonnet 给出 $\chi>0$，由分类公式可知它拓扑上为球面。这不说明它必须是几何上的圆球，椭球已是反例。若处处 $K=0$，公式仅给出 $\chi=0$；抽象平坦环面确实存在，但这一点不意味着它能作为 $\R^3$ 中光滑正则曲面等距嵌入。

非紧曲面不能直接套用无边界整体公式。例如平面 $K=0$，却没有 $0=2\pi\chi(\R^2)$ 这样的结论；若在平面中取一个圆盘，边界项恰为 $2\pi$。带边界的区域若漏掉内边界方向，也会得出错误结果：平面圆环的外圆正向为逆时针，内圆正向为顺时针，两条边界曲率积分分别为 $2\pi$ 和 $-2\pi$，总和为零，与 $\chi=0$ 一致。

平行移动同样反映曲率。向量相对局部正交标架的角度为 $\psi$ 时，平行条件为 $\psi'+\omega(\dot\gamma)=0$。沿坐标圆盘边界闭合移动，向量最终相对初始方向的旋转角模 $2\pi$ 为 $-\int_{\partial D}\omega=\int_DK\dd A$。因此小回路的平行移动效应由所围区域的高斯曲率积分控制。

\section{习题}
\begin{exercise}\label{ex:dg-gb-triangle}
半径 $R$ 球面上一个由短大圆弧围成的三角形的三个内角均为 $2\pi/3$，求其面积。
\end{exercise}
\begin{exercise}\label{ex:dg-gb-annulus}
用边界 Gauss--Bonnet 公式核对平面圆环 $a\le\sqrt{x^2+y^2}\le b$，其中 $0<a<b$，特别说明两个边界方向。
\end{exercise}
\begin{exercise}\label{ex:dg-gb-positive}
设闭合、连通、可定向曲面满足 $K>0$。证明其亏格为零，并说明为何不能由此推出曲面一定是标准球面。
\end{exercise}
\begin{exercise}\label{ex:dg-gb-flat-embedded}
证明不存在紧致无边界正则曲面 $\Sigma\subset\R^3$ 满足 $K\le0$ 处处成立。提示：考察到固定点的平方距离在曲面上的最大值。
\end{exercise}

